\section{Introduction}
\label{sec:introduction}

The 2019 decision in \textit{Rucho v.\ Common Cause}\footnote{\textit{Rucho v.\ Common Cause}, 588 U.S. 684 (2019).} produced a paradox of democratic governance. The Supreme Court held that federal courts lack judicially manageable standards for adjudicating partisan gerrymandering claims, effectively withdrawing the federal judiciary from oversight of what Chief Justice Roberts himself acknowledged was a ``unjust'' and ``incompatible with democratic principles'' practice.\footnote{\textit{Id.} at 693.} The Court directed aggrieved voters to seek relief through state courts interpreting state constitutions, and through federal legislation. Both pathways remain largely unexplored.

This Article takes up the legislative pathway. The central claim is that algorithmic redistricting---the use of computer programs that draw district boundaries by optimizing population equality and compactness without accessing any partisan data---is legally viable, constitutionally sound, and ready for legislative adoption at both the federal and state levels. The argument proceeds through three adoption pathways: a federal statute enacted pursuant to the Elections Clause of Article I, Section 4; state legislation and commission action; and court-ordered remedial use. For each pathway, the Article identifies the constitutional authority, analyzes the relevant doctrine, and provides model text.

The timeliness of this analysis is not merely academic. The 2020 redistricting cycle produced congressional maps that federal courts largely declined to invalidate on partisan grounds, following \textit{Rucho}. State courts in Pennsylvania, North Carolina, and New York have been more active, but their involvement is episodic and their remedies inconsistent.\footnote{See \textit{League of Women Voters v.\ Commonwealth}, 178 A.3d 737 (Pa. 2018); \textit{Harper v.\ Hall}, 868 S.E.2d 499 (N.C. 2022); \textit{Harkenrider v.\ Hochul}, 197 A.D.3d 88 (N.Y. App. Div. 2021).} The underlying problem---legislatures drawing districts to entrench their own majorities---has not been solved by any institutional mechanism currently in place. Algorithmic redistricting offers a structural solution: if the process cannot be manipulated because the mapmaker cannot see partisan data, the problem of partisan manipulation is not merely ameliorated but eliminated.

The Article proceeds as follows. Section~\ref{sec:constitutional} surveys the constitutional foundation for federal and state redistricting authority, centering on the Elections Clause and its interpretation in \textit{Smiley v.\ Holm} and \textit{Arizona State Legislature v.\ Arizona Independent Redistricting Commission}. Section~\ref{sec:federal} develops the federal statutory pathway, with the Huntington-Hill apportionment statute (2 U.S.C. \S\,2a) as a structural precedent for mathematical mandates in the apportionment-redistricting domain. Section~\ref{sec:state} analyzes state adoption pathways, categorizing states by constitutional compatibility. Section~\ref{sec:vra} addresses Voting Rights Act compliance, arguing that the algorithm's VRA mode satisfies the \textit{Gingles} prongs without partisan contamination and creates a safe harbor for majority-minority districts. Section~\ref{sec:challenges} analyzes anticipated legal challenges---nondelegation, Elections Clause limits, and VRA conflict---and explains why each fails. Section~\ref{sec:model} presents model statute text in U.S. Code drafting style. Section~\ref{sec:conclusion} concludes.

\subsection{The Problem Algorithmic Redistricting Solves}

Before analyzing the law, it is worth specifying precisely what problem algorithmic redistricting solves. The gerrymandering problem has two distinct components that are often conflated: the \textit{intent problem} and the \textit{structural problem}.

The intent problem is that redistricting bodies draw maps with partisan advantage in mind. Courts have struggled to police this because partisan intent is difficult to prove, and because the line between legitimate partisan considerations and unlawful partisan excess has proven impossibly difficult to define. \textit{Rucho} represents the federal judiciary's capitulation to this difficulty.

The structural problem is prior to intent: redistricting bodies have access to, and routinely use, fine-grained partisan data at the census tract and block level. This data access is what makes precise gerrymandering possible. Modern redistricting software can display, in real time, the effect of every boundary adjustment on each party's expected seat count. Even a well-intentioned mapmaker working alone with this software would find it difficult to ignore the partisan implications of their choices.

Algorithmic redistricting solves both problems simultaneously, but in fundamentally different ways. It solves the structural problem by design: the algorithm receives only population counts and geographic adjacency data from the census, and is architecturally incapable of accessing voter registration, election results, or any partisan information. It cannot gerrymander not because its operators are virtuous but because it lacks the inputs that gerrymandering requires. It solves the intent problem derivatively: because the process has no access to partisan data, no partisan intent can be operationalized through it. The algorithm's outputs may still correlate with partisan advantage---because where people live is correlated with how they vote---but the correlation reflects geography, not manipulation.\footnote{See \citet{dellaLibera2026recursive} (demonstrating that geographic sorting produces a Democratic advantage in algorithmic maps that mirrors the efficiency gap literature's findings, but tracing this advantage to the urban-rural sorting of the electorate rather than any strategic design).}

This distinction matters for the legal analysis. Arguments that algorithmic redistricting violates the Constitution or the VRA typically assume that algorithmic outputs should be evaluated against partisan outcomes as standards. But the correct standard for evaluating an algorithm that cannot see partisan data is not partisan outcome---it is the reasonableness of the criteria it does optimize. On this standard, the case is straightforward: population equality and geographic compactness are unambiguously legitimate redistricting criteria that every court and state constitution endorses.
